\section{Experimental Protocol}
\label{sec:setup}

\subsection{Data, optimization, and run identity}

We analyze an ImageNet-1k classification campaign with seed 42. Evaluation
uses the 50,000-image validation split and reports top-1 accuracy of the
exponential moving average (EMA) model. Every available training summary row
was matched to its final per-epoch EMA validation line in the launcher log:
2,419 rows across 17 runs with recorded epochs. There are 18 launched runs in
total: 16 complete schedules, one partial curriculum snapshot, and one
512-pixel \mn{} run with no completed epoch in the supplied record.
Checkpoint-verification receipts are present, but the weights themselves
were not transferred with the campaign packet. This audit verifies the
recorded outcomes; it is not an independent checkpoint re-evaluation.

The 224-pixel runs train from scratch for 150 or 300 epochs, using AdamW,
cosine learning-rate decay, five warmup epochs (twenty for the warmup probe),
weight decay 0.05, gradient
clipping at 1.0, label smoothing 0.1, Mixup 0.8, CutMix 1.0, RandAugment, and
random erasing probability 0.25. The base learning rate is $5\times10^{-4}$;
a separate \mn{} arm uses $7.5\times10^{-4}$. The higher-rate arm supplies
the 300-epoch \mn{} checkpoint. Therefore the final dense-versus-\mn{}
comparison is a comparison of selected configurations, while the geometry
comparison uses matched learning rates. No teacher distillation is enabled.
Appendix~\ref{app:recipe} gives the recipe and auxiliary-logit schedule.

The effective global batch is 1,024 throughout the recorded campaign. The
384/512-pixel fine-tunes use four-step gradient accumulation:
$32\times8\times4=1024$ for the dense 384-pixel run and
$64\times4\times4=1024$ for the other fine-tunes. Each fine-tune epoch contains
5,004 per-rank microsteps and 1,251 optimizer steps. Counting microsteps as
optimizer steps would incorrectly suggest a fourfold batch reduction.
Different world sizes and microbatches still produce different execution
and sample-grouping conditions; effective batch equality alone does not
establish bitwise equivalence.

\subsection{Resolution adaptation and checkpoint selection}

Three 384-pixel runs fine-tune the corresponding 300-epoch EMA weights for
30 epochs, using learning rate $10^{-5}$, minimum rate $10^{-6}$, one warmup
epoch, weight decay $10^{-8}$, and no Mixup or CutMix. EMA decay changes from
0.99996 to 0.9996. The \mn{} arms either retain $R=3$ or use
$R=3\times384/224\approx5.143$. The latter tests a particular proportional
radius scaling; it is not a complete radius search. Bicubic validation
preprocessing uses crop percentage 0.9 at 224 pixels and 1.0 for the
fine-tunes. The completed 512-pixel dense result follows its 384-pixel
checkpoint; there is no completed paired \mn{} result at that resolution.

We report both the best observed EMA top-1 and the final-epoch EMA top-1.
Best checkpoints are selected on the same validation split used for
reporting, so these scores are not independent held-out estimates of model
selection. Epoch indices start at zero. All completed-run comparisons use
actual summary headers rather than fixed column positions, because the
summary schema changes with enabled diagnostics. Runs with a teardown error
after a complete recorded schedule are distinguished from unfinished runs.

\subsection{Post-training compression and timing}

The accuracy sweep applies ToMe and PiToMe to the dense EMA checkpoint
without additional training, and changes \mn{}'s carrier budget without
retraining each point. Thus \mn{} has been trained with its native
compression, whereas the dense backbone was trained without merging. This
is a deployment comparison of available checkpoints, not an isolation of the
routing rule under identical training. We distinguish a 224-pixel checkpoint
evaluated at a new resolution from a checkpoint fine-tuned at that resolution.
Repeated successful sweep records are resolved by latest timestamp, not
highest accuracy. All included accuracy cells use the full validation split. Training-summary
validation uses 64 images per rank. For sweeps from the 224-pixel checkpoints,
\mn{} and ToMe use batch 128 and PiToMe uses batch 256; the final 384-pixel
fine-tuned curves use batch 256 for all three. The separate fine-tuned gate
uses batch 384. These batch distinctions matter for the soft-routing
implementation and are retained in the evidence ledger.

The 224-pixel efficiency experiment uses one H20-3e, synthetic inputs, batch
64, fp16 autocast, 30 warmup iterations and 50 measured iterations. It records
CUDA-event medians and interquartile ranges (IQRs); training cells include
forward, loss, backward and AdamW update. Device isolation is checked around
the measurements. The separate radius-3 benchmarks at 384 pixels use the
same thirty-warmup/fifty-measurement protocol at batches 64 and 128.
These are model-execution timings, excluding image loading,
and do not estimate end-to-end training time on eight GPUs. Training-log
throughput is excluded because the training machine was shared.

A later 384-pixel benchmark uses batch 64, ten warmup and 30 measured
iterations, random model initialization and synthetic inputs. It records
achieved token counts, but disables ToMe proportional attention, whereas the
accuracy sweep uses it. We retain these latency measurements as a separate
configuration-specific experiment and do not attach accuracy coordinates to
them. Final token counts also require care: the accuracy sweep counts patches,
while this timing harness counts the class token as well.
